\begin{helper}
For the natural reduced-vertex count and \(\beta=1/2\), injections into the
two-coordinate reduced space exist eventually, and the total
\(\noderef{TwoBiteEventProbability}\) mass of the whole two-bite sample space
is then exactly \(1\).
\[
\forall^{\mathrm{eventually}} n,\qquad
n\le m(n)^2
\quad\text{and}\quad
\noderef{TwoBiteEventProbability}
  \bigl(n,m(n),\noderef{TwoBiteEdgeProbability}(\beta,n),\mathrm{True}\bigr)=1.
\]
\end{helper}
\begin{proof}
Write \(L=\log n\) and
\[
M=m(n)=\noderef{TwoBiteNaturalReducedVertexCount}(n)
  =\left\lfloor \frac{n}{L^2}\right\rfloor
\]
for all \(n\), using the hypothesis on \(m\).  The equality
\(\beta=1/2\) only identifies the edge parameter; the normalization of the
whole sample space works for any real value of that parameter.

First prove eventual injection support.  Since \(L\to\infty\) and
\(L^4=o(n)\), for all sufficiently large \(n\) we have \(L>0\),
\[
\frac{n}{L^2}\ge2,\qquad 4L^4\le n .
\]
For \(x\ge2\), \(\lfloor x\rfloor\ge x-1\ge x/2\).  Applying this to
\(x=n/L^2\) gives
\[
M\ge \frac{n}{2L^2}.
\]
Therefore
\[
M^2\ge \frac{n^2}{4L^4}\ge n,
\]
where the last inequality is exactly \(4L^4\le n\).  Hence
\[
n\le M^2=|\operatorname{Fin}M\times\operatorname{Fin}M|.
\]
For finite types, an embedding \(\operatorname{Fin} n\hookrightarrow
\operatorname{Fin}M\times\operatorname{Fin}M\) exists whenever the target
cardinality is at least the source cardinality.  Equivalently, the type
\(\operatorname{Fin} n\hookrightarrow\operatorname{Fin}M\times\operatorname{Fin}M\)
has positive cardinality.

Now expand the total mass.  Let
\[
C_n=\operatorname{Fin} n\hookrightarrow(\operatorname{Fin}M\times
\operatorname{Fin}M)
\]
and set \(p=\noderef{TwoBiteEdgeProbability}(\beta,n)\).  Unfolding
\(\noderef{TwoBiteEventProbability}\) and
\(\noderef{TwoBiteSampleWeight}\), the total mass of the event
\(\mathrm{True}\) is
\[
\sum_{G_R,G_B,\pi}
\noderef{GnpGraphWeight}(p,G_R)\,
\noderef{GnpGraphWeight}(p,G_B)\,
\noderef{UniformInjectionWeight}(n,M).
\]
This finite sum factors as the product of the red graph total mass, the blue
graph total mass, and the injection-coordinate total mass.  The two graph
factors are each \(1\) by \(\noderef{GnpGraphWeightSumOne}\).

Because \(C_n\) has positive cardinality, the reciprocal branch in
\(\noderef{UniformInjectionWeight}\) applies:
\[
\noderef{UniformInjectionWeight}(n,M)=|C_n|^{-1}.
\]
Thus the injection-coordinate sum is
\[
\sum_{\pi\in C_n}|C_n|^{-1}
  = |C_n|\cdot |C_n|^{-1}=1.
\]
Multiplying the three factors gives total mass
\[
1\cdot1\cdot1=1.
\]
The construction holds for all sufficiently large \(n\), so both the
injection-support inequality and the exact normalization statement hold
eventually.
\end{proof}
